\subsection{Equações do Movimento de Atitude da Espaçonave Visitante no Referencial Inercial}
\label{sec:atitude_inercial}

A formulação matemática de atitude da espaçonave visitante é realizada partir das equações de Euler aplicadas a um corpo rígido.
Essas equações são descritas em torno do centro de massa (\emph{Center of Mass} ou CoM) e expressas no sistema da própria espaçonave.
As equações diferenciais de movimento rotacional assumem a forma:

\begin{equation}
\mathbf{I}\,\dot{\vec{\omega}}
+ \vec{\omega}\times \big(\mathbf{I}\,\vec{\omega}\big)
= \vec{M},
\label{eq:euler_corpo_rigido}
\end{equation}

em que $\mathbf{I}$ é a matriz de inércia em relação ao centro de massa, $\vec{\omega}$ é o vetor de velocidade angular e $\vec{M}$ são os momentos aplicados ao corpo.

\subsubsection{Matriz de inércia variável}

No caso da espaçonave visitante, que possui um MR acoplado em sua base, a matriz de inércia equivalente não é constante, pois a movimentação dos elos modifica a distribuição de massa, sendo assim, a inércia total do sistema pode ser representada por:

\begin{equation}
\mathbf{I}(t) = \mathbf{I}_{base} + \mathbf{I}_{manip}(t),
\label{eq:inercia_total}
\end{equation}

em que $\mathbf{I}_{base}$ é a matriz associada à base da espaçonave e $\mathbf{I}_{manip}(t)$ é a contribuição variável do manipulador robótico.

\subsubsection{Dinâmica acoplada com o manipulador}

Os movimentos do MR geram torques que devem ser incluídos na equação dinâmica da atitude, \eqref{eq:euler_corpo_rigido}.
Deve-se considerar tanto os torques de controle externos quanto os torques internos de reação do MR, dada por:

\begin{equation}
\mathbf{I}(t)\,\dot{\vec{\omega}}
+ \vec{\omega}\times \big(\mathbf{I}(t)\,\vec{\omega}\big) 
= \vec{M}_{controle} + \vec{M}_{manip},
\label{eq:dinamica_acoplada}
\end{equation}

onde $\vec{M}_{controle}$ são os torques de controle aplicados à espaçonave e $\vec{M}_{manip}$ são os torques produzidos pelas juntas do MR durante sua movimentação. 
